\documentclass[10pt]{article}

\usepackage[a4paper,margin=22mm]{geometry}
\usepackage[T1]{fontenc}
\usepackage{lmodern}
\usepackage{microtype}
\usepackage{booktabs}
\usepackage{xcolor}
\usepackage{titlesec}
\usepackage{caption}
\usepackage{enumitem}
\usepackage{array}
\usepackage{fancyvrb}
\usepackage[hidelinks]{hyperref}

\definecolor{secblue}{HTML}{1F4E8C}
\definecolor{rulegrey}{HTML}{888888}
\definecolor{codebg}{HTML}{F1F5F9}

\titleformat{\part}{\Large\bfseries\color{secblue}}{Part \thepart.}{0.6em}{}
\titleformat{\section}{\large\bfseries\color{secblue}}{\thesection}{0.6em}{}
\titleformat{\subsection}{\normalsize\bfseries\color{secblue}}{\thesubsection}{0.6em}{}
\titlespacing*{\section}{0pt}{12pt}{5pt}
\titlespacing*{\subsection}{0pt}{8pt}{3pt}

\captionsetup{font=small,labelfont=bf,labelsep=period}
\setlist[itemize]{leftmargin=1.4em,itemsep=2pt,topsep=3pt}
\setlist[enumerate]{leftmargin=1.6em,itemsep=2pt,topsep=3pt}
\renewcommand{\arraystretch}{1.15}

% inline code
\newcommand{\code}[1]{\texttt{\small #1}}

% verbatim diagram block on a light background
\DefineVerbatimEnvironment{diag}{Verbatim}{fontsize=\scriptsize,frame=single,
  framerule=0.3pt,rulecolor=\color{rulegrey},xleftmargin=4pt,xrightmargin=4pt}
\DefineVerbatimEnvironment{ccode}{Verbatim}{fontsize=\footnotesize,frame=leftline,
  framerule=1.2pt,rulecolor=\color{secblue},xleftmargin=8pt}

\setlength{\parindent}{0pt}
\setlength{\parskip}{5pt}

\begin{document}

% ---------------------------------------------------------------- title block
\begin{center}
  {\LARGE\bfseries cMCP --- Study Guide}\\[4pt]
  {\normalsize The Model Context Protocol, and how a from-scratch C11 library implements it.\\
   A document to \emph{learn} the protocol and to \emph{read} this codebase.}\\[6pt]
  {\small\color{rulegrey} Generated 2026-06-09 \quad\textbullet\quad
   cMCP \code{v0.9.0} \quad\textbullet\quad protocol \code{2025-11-25} \quad\textbullet\quad pure C11}
\end{center}
\vspace{-2pt}
{\color{secblue}\hrule height 1.2pt}
\vspace{6pt}

{\small\textbf{How to read this.} Part~I teaches the Model Context Protocol from
first principles --- no C required, MCP as a specification. Part~II walks the
cMCP source tree layer by layer and shows \emph{where} each protocol concept is
implemented. Part~III is a guided reading order plus hands-on exercises. If you
already know MCP, skim Part~I for cMCP's vocabulary and jump to Part~II. The
companion documents are \code{docs/architecture.md} (the ``why'' of the
partitioning, in prose) and \code{docs/cmcp-engineering-report.tex} (the
\emph{measured} conformance / performance / hardening claims). This guide is the
\emph{didactic} layer over both.}

\tableofcontents
\vspace{4pt}
{\color{rulegrey}\hrule height 0.6pt}

% =====================================================================
\part{The Model Context Protocol}

% ---------------------------------------------------------------------
\section{What MCP is, and the problem it solves}

A large language model is a function from text to text. To do anything in the
world --- read a file, search a corpus, call an API --- it needs an external
program to act on its behalf and feed results back into its context window. Before
MCP, every host application (an IDE, a chat client, an agent) invented its own
bespoke plugin protocol, and every tool author re-implemented against each host.
That is an $M\times N$ integration problem.

The \textbf{Model Context Protocol} (introduced by Anthropic as an open standard)
collapses it to $M+N$: tools speak one protocol, hosts speak one protocol, and
any host can drive any tool. The common analogy is ``a USB-C port for AI
applications'' --- a single physical shape that many devices plug into.

\textbf{Three roles.} An MCP deployment has exactly three kinds of participant:

\begin{itemize}
  \item \textbf{Host} --- the application the user interacts with, and the thing
    that owns the LLM (e.g.\ Claude Desktop, an IDE, or cMCP's target consumer,
    the \emph{butlerbot} agent). The host decides which servers to connect to and
    whether to honour a server's requests.
  \item \textbf{Client} --- the connector \emph{inside} the host that maintains a
    \emph{one-to-one} session with a single server. A host with three servers
    runs three clients. In cMCP this is \code{libcmcp\_client.a}.
  \item \textbf{Server} --- a separate program that exposes capabilities (tools,
    resources, prompts) over the protocol. In cMCP this is
    \code{libcmcp\_server.a}; the reference servers are \code{filesystem-mcp} and
    \code{crag-mcp}.
\end{itemize}

\begin{diag}
   +------------------------- Host application -------------------------+
   |   LLM   <-->   client A   <-->   client B   <-->   client C        |
   +------------------|---------------|----------------|----------------+
                      |               |                |  (one client per server,
                      v               v                v   a 1:1 session each)
                  Server A        Server B         Server C
                (filesystem)     (retrieval)       (custom)
\end{diag}

The key mental model: the \textbf{client/server link is symmetric in machinery
but asymmetric in role}. Both sides frame the same JSON-RPC, negotiate
capabilities, and can send unsolicited messages --- but the server \emph{offers}
primitives and the client \emph{drives} them. cMCP exploits this directly: $\sim$60\%
of the code (framing, schema, transport, capability structs) is shared in
\code{libcmcp\_core.a}, and only the asymmetric halves live in the server and
client libraries.

% ---------------------------------------------------------------------
\section{The JSON-RPC 2.0 foundation}

Every MCP message is a JSON-RPC 2.0 object. There are exactly three shapes, and
understanding them is most of understanding the wire.

\begin{enumerate}
  \item \textbf{Request} --- carries \code{"id"}, \code{"method"}, optional
    \code{"params"}. Demands a response. \emph{Example:}
    \code{\{"jsonrpc":"2.0","id":7,"method":"tools/call","params":\{\dots\}\}}
  \item \textbf{Response} --- carries the same \code{"id"} and \emph{either}
    \code{"result"} (success) \emph{or} \code{"error"} (never both).
  \item \textbf{Notification} --- a request with \emph{no} \code{"id"}. Fire and
    forget; no response is ever sent. \emph{Example:}
    \code{notifications/initialized}, \code{notifications/cancelled}.
\end{enumerate}

\textbf{The id space.} Ids match responses to requests. Each session has a
single monotonic id space. The side that \emph{sends} a request owns the id and
correlates the inbound response by it. Because both sides may send requests
(the server can ask the host for an LLM completion, see \S5), there are
\emph{two} id spaces per session --- one per direction --- and they are
independent.

\textbf{Error codes.} JSON-RPC defines a standard set, all reused by MCP:

\begin{center}\small
\begin{tabular}{r l l}
\toprule
Code & Name & Meaning \\
\midrule
\code{-32700} & Parse error      & Invalid JSON received \\
\code{-32600} & Invalid request  & Not a valid Request object (e.g.\ a batch, which MCP forbids) \\
\code{-32601} & Method not found & Unknown method, or a default-deny decline \\
\code{-32602} & Invalid params   & Bad/missing params (e.g.\ malformed \code{protocolVersion}) \\
\code{-32603} & Internal error   & Server-side failure (e.g.\ structured output fails its schema) \\
\bottomrule
\end{tabular}
\end{center}

\textbf{No batching.} JSON-RPC 2.0 permits arrays of messages; MCP \emph{removed}
batch support (since spec \code{2025-06-18}). A conformant implementation parses a
batch only so it can reject it cleanly with \code{-32600}.

\textbf{The dual error-channel subtlety (load-bearing).} There are \emph{two}
distinct ways a \code{tools/call} can ``fail,'' and a correct host must check
both:
\begin{itemize}
  \item A \textbf{protocol-level} failure (unknown tool, malformed params) is a
    JSON-RPC \code{error} object with code \code{-32602}/\code{-32601}.
  \item A \textbf{tool-level} failure (the tool ran but its arguments were
    invalid, or it errored semantically) is a \emph{successful} JSON-RPC
    \code{result} whose body is
    \code{\{isError:true, content:[\{text:"\dots"\}]\}}.
\end{itemize}
This split is deliberate in the spec: ``the tool could not be invoked'' is a
protocol error; ``the tool ran and reports a problem'' is data the model should
see. cMCP matches the TypeScript reference SDK here and pins it by test (\S13).

% ---------------------------------------------------------------------
\section{Lifecycle: handshake and capability negotiation}

An MCP session has a strict opening sequence. Nothing operational happens until
it completes.

\begin{diag}
   client                                                          server
   ------                                                          ------
   initialize (protocolVersion, clientInfo, capabilities) ------->
                <------ result (protocolVersion, serverInfo, capabilities)
   notifications/initialized ------------------------------------>
   =================== session is now "ready"; normal traffic ===============
\end{diag}

\textbf{Step 1 --- \code{initialize}.} The client sends its protocol version, an
\code{clientInfo} \code{\{name, version\}}, and a \emph{capabilities} object
declaring what \emph{it} can do (sampling? roots? elicitation?). The server
replies with its own version, \code{serverInfo}, and \emph{its} capabilities
(tools? resources? prompts? logging?).

\textbf{Step 2 --- \code{notifications/initialized}.} The client confirms it
processed the result. Only now may either side issue operate-class requests. A
server that receives, say, \code{tools/call} before \code{initialized} answers
\code{-32600}.

\textbf{Version tolerance (a common mistake).} The protocol version is a date
string (cMCP pins \code{2025-11-25}). If the peer advertises a \emph{different}
version, this is \textbf{not} a handshake failure: the spec says respond with
success, advertise your own version, and let the \emph{other} side decide whether
to disconnect. Only a \emph{missing or malformed} \code{protocolVersion} is a hard
\code{-32602}. cMCP captures the peer's version (\code{cmcp\_client\_server\_protocol})
and proceeds.

\textbf{Capability negotiation is the contract.} A capability not advertised must
not be used. If a server never declared \code{resources.subscribe}, the client
must not call \code{resources/subscribe}; if a client never declared
\code{sampling}, the server must not ask it for a completion. This is the
mechanism that keeps both ends from sending messages the other isn't prepared to
handle. Sub-capabilities (\code{listChanged}, \code{subscribe}, \code{logging})
are opt-in flags nested under the top-level capability.

% ---------------------------------------------------------------------
\section{The three server primitives}

A server exposes three kinds of thing. Each has a \code{*/list} method to
enumerate it and a method to act on one. Knowing who controls each is the
clarifying lens (the spec frames this as \emph{who is in the driver's seat}):

\begin{center}\small
\begin{tabular}{l l l l}
\toprule
Primitive & Control model & List & Act \\
\midrule
\textbf{Tools}     & model-controlled    & \code{tools/list}     & \code{tools/call} \\
\textbf{Resources} & application-controlled & \code{resources/list} & \code{resources/read} \\
\textbf{Prompts}   & user-controlled     & \code{prompts/list}   & \code{prompts/get} \\
\bottomrule
\end{tabular}
\end{center}

\textbf{Tools} are functions the \emph{model} chooses to invoke. Each carries a
\code{name}, a human \code{description}, and an \code{inputSchema} (JSON Schema)
describing its arguments. The model reads the schema, fills arguments, and the
host issues \code{tools/call}. Results are a \code{content} array of typed items
(\code{text}, \code{image}, \code{resource\_link}, \dots), optionally a typed
\code{structuredContent} validated against an \code{outputSchema}.

\textbf{Resources} are data the \emph{application} pulls into context --- a file,
a database row, a diagnostics blob. Each has a \code{uri} (e.g.\
\code{file:///path}, \code{crag://stats}). The host reads them with
\code{resources/read}. A client may \code{resources/subscribe} to a URI and
receive \code{notifications/resources/updated} when it changes (only if the
server declared \code{resources.subscribe}).

\textbf{Prompts} are reusable, parameterised message templates the \emph{user}
invokes (think slash-commands). \code{prompts/list} enumerates them with a flat
argument descriptor list \code{[\{name, description?, required?\}]} --- note this
is \emph{not} full JSON Schema --- and \code{prompts/get} expands one with
arguments into a message list.

\textbf{list\_changed.} If a server's primitive set is dynamic, it declares
\code{listChanged} and emits \code{notifications/tools/list\_changed} (etc.) so
the client knows to re-list.

% ---------------------------------------------------------------------
\section{The reverse direction: server-initiated requests}

MCP is bidirectional. Beyond responding to the client, a server can originate
requests \emph{back to the host} --- this is what makes MCP an agent protocol
rather than an RPC catalogue. There are three such directions, all gated by a
\emph{client} capability:

\begin{itemize}
  \item \textbf{Sampling} (\code{sampling/createMessage}) --- the server asks the
    host's LLM to produce a completion. Useful for a tool that wants to summarise
    or re-reason over raw output before returning it. Gated by the client's
    \code{sampling} capability; default-decline (\code{-32601}) if the host
    registered no handler.
  \item \textbf{Roots} (\code{roots/list}) --- the server asks the host which
    filesystem paths / URIs are in scope. A filesystem-shaped server consults
    this before touching anything. The host declares \code{roots} and may emit
    \code{notifications/roots/list\_changed}.
  \item \textbf{Elicitation} (\code{elicitation/create}) --- mid-tool-call, the
    server asks the user for structured input: a confirmation, a missing
    argument, a credential. The request carries a \code{requestedSchema}; the
    host returns \code{\{action:"accept"|"decline"|"cancel", content\}}.
\end{itemize}

Two more cross-cutting mechanisms ride alongside:

\begin{itemize}
  \item \textbf{Progress} --- a long call attaches a \code{progressToken} (in
    \code{params.\_meta}); the server emits \code{notifications/progress} carrying
    that token, and the host correlates it to the originating call.
  \item \textbf{Cancellation} --- either side may send
    \code{notifications/cancelled \{requestId, reason?\}} to abandon an in-flight
    request. It is \emph{cooperative}: the receiver should stop work, but a
    response that crosses the cancel on the wire is simply dropped.
  \item \textbf{Logging} --- with the \code{logging} capability, a server ships
    structured \code{notifications/message \{level, logger?, data\}} events, and
    the host dials the floor with \code{logging/setLevel}. Levels are the eight
    RFC 5424 syslog levels.
\end{itemize}

\textbf{The trust model.} Every reverse request is default-deny. The host opts in
\emph{per server} by registering a handler --- ``did the host bother to register a
handler for this server?'' is the authorisation gate. There is no global
allow-list; trust is decided one server at a time.

% ---------------------------------------------------------------------
\section{Transports}

MCP separates message \emph{semantics} from message \emph{transport}. The spec
defines two standard transports; cMCP implements both.

\begin{itemize}
  \item \textbf{stdio} --- the server is a subprocess; the host writes JSON-RPC to
    its stdin and reads from its stdout, one message per line (newline-delimited).
    \emph{stdout is sacred} --- any stray write corrupts the wire, so all logging
    goes to stderr. This is the primary transport for local, single-tenant tools
    and cMCP's design centre.
  \item \textbf{Streamable HTTP} --- a single \code{/mcp} endpoint. The client
    \code{POST}s a request and gets the response in the body; it may also open a
    \code{GET} with \code{Accept: text/event-stream} to receive server-initiated
    messages over Server-Sent Events (SSE). A session id (\code{Mcp-Session-Id}
    header) is minted on the first \code{initialize} and threaded through
    subsequent calls. SSE events carry monotonic ids so a dropped stream can
    resume via \code{Last-Event-Id} (spec SEP-1699).
\end{itemize}

\textbf{Why stdio is the default for agents.} A local tool launched as a
subprocess needs no ports, no auth, no TLS --- the OS process boundary \emph{is}
the security boundary, and stdin/stdout are already private to the parent. HTTP
matters for remote/shared servers but brings a network attack surface (\S12).

% =====================================================================
\part{Inside cMCP}

% ---------------------------------------------------------------------
\section{Build, link targets, and the six-layer pipeline}

cMCP is one source tree compiled into three static libraries plus reference
binaries. The partitioning mirrors the protocol's shared-machinery /
asymmetric-role split from \S1.

\begin{diag}
  libcmcp_core.a    json + rpc + schema + types + transports        (shared ~60%)
  libcmcp_server.a  + server.c + worker.c        (depends on core)   (server half)
  libcmcp_client.a  + client.c + session.c       (depends on core)   (client half)
\end{diag}

Link order matters: consumers (server/client) before provider (core), so the
linker resolves core symbols on the second pass. The Makefile handles it. The
core libraries are dependency-free; only the HTTP \emph{client} links
\code{libcurl}.

\textbf{Build commands to know:}
\begin{center}\small
\begin{tabular}{l l}
\toprule
Command & Effect \\
\midrule
\code{make}            & libs + \code{cmcp-inspect} + \code{filesystem-mcp} + \code{cmcp-tee} + examples \\
\code{make test}       & build and run every \code{tests/test\_*.c} (offline, hermetic) \\
\code{make valgrind}   & same suite under Valgrind (leak-free) \\
\code{make test-asan}  & rebuild under ASan/UBSan, run suite \\
\code{make test-tsan}  & rebuild under ThreadSanitizer, run suite \\
\code{make crag-mcp}   & build the cRAG reference server (needs sibling \code{../cRAG}) \\
\code{make conformance}& cross-check vs the TS reference SDK (needs Node + network) \\
\bottomrule
\end{tabular}
\end{center}

The request pipeline is six layers, each in its own module, each talking only to
the layer directly below:

\begin{diag}
  +----------------------------------------------------------------+
  |                      tools (handlers)                          |  <- user code
  +----------------------------------------------------------------+
  |  server.c - registry, dispatch    |  client.c - async client   |
  |   + worker.c pool                 |   + session.c N-server mux  |
  +----------------------------------------------------------------+
  |  schema.c - JSON Schema subset (validates tools/call.arguments)|
  +----------------------------------------------------------------+
  |  rpc.c - JSON-RPC 2.0 framing, in-flight ID table              |
  +----------------------------------------------------------------+
  |  transport_stdio.c / transport_http.c / transport_http_client.c|
  +----------------------------------------------------------------+
  |  json.c - typed value tree, parser, emitter                    |
  +----------------------------------------------------------------+
\end{diag}

The discipline is strict: the transport layer never inspects message content; the
RPC layer never reads raw bytes; the schema layer never speaks JSON-RPC. That
clean separation is what lets the same RPC core serve both halves and lets HTTP
plug in on the same vtable as stdio.

% ---------------------------------------------------------------------
\section{Layer by layer}

\subsection{json.c --- the value tree (\S6, layer 1)}

A hand-written JSON parser/emitter over a discriminated-union value tree
(\code{cmcp\_json\_t}). No third-party JSON library. Lifted from cRAG's
\code{util.c} and extended with object/array constructors, escape-aware string
emission (UTF-8 surrogate pairs, control bytes), and \emph{stable key ordering on
emit} --- keys come out in insertion order. That last is not a JSON requirement
but it makes wire frames byte-stable, which is what lets tests and the replay gate
diff frames literally.

Hardening lives here too: \code{CMCP\_JSON\_MAX\_DEPTH} (default 64) bounds
recursion against stack exhaustion, and \code{CMCP\_JSON\_MAX\_ELEMENTS} (default
65536) caps array/object size against memory bombs.

\subsection{rpc.c --- JSON-RPC framing (layer 2)}

A discriminated-union message type (\code{cmcp\_rpc\_message\_t}) covering
request / response / notification, plus emit and parse. It parses batches only to
reject them with \code{-32600} (\S2). The \emph{in-flight ID table}
(\code{cmcp\_rpc\_pending\_t}) is the single monotonic id space: the client
reserves an id before sending, the reader looks the id up on the inbound response
and routes the parsed message to whoever waits. \code{CMCP\_RPC\_MAX\_INFLIGHT}
(default 1024) caps outstanding ids.

\subsection{schema.c --- the JSON Schema subset (layer 3)}

The server validates inbound \code{tools/call.arguments} against the tool's
registered \code{inputSchema} \emph{before} the handler runs --- so handlers only
ever see well-formed input. The same validator runs \emph{outbound} when a handler
attaches \code{structuredContent} against an \code{outputSchema} (mismatch ->
\code{-32603}).

Supported keywords: \code{type} (incl.\ array-of-types), \code{properties},
\code{required}, \code{additionalProperties:false}, \code{enum},
\code{minLength}/\code{maxLength} (Unicode code points), \code{minimum}/\code{maximum},
\code{items}. \emph{Deliberately unsupported} in v0.1: \code{\$ref},
\code{oneOf}/\code{anyOf}/\code{allOf}, \code{format}, \code{pattern}. This is a
documented subset (\code{docs/schema-conformance.md}), not a silent gap --- those
keywords would multiply validator complexity for cases real MCP tools don't need.

Recall the dual channel from \S2: a \code{tools/call} \emph{argument} failure
surfaces as a tool-level \code{\{isError:true\}} result, while protocol-level
rejections use the \code{-32602} error channel.

\subsection{transport\_*.c --- the vtable (layer 4)}

All three transports implement one vtable, \code{cmcp\_transport\_t}:
\code{read\_fn}, \code{write\_fn}, \code{close\_fn}, and an optional
\code{wake\_fn}.

\textbf{stdio} (\code{transport\_stdio.c}) --- newline-delimited JSON over a
\code{(read\_fd, write\_fd)} pair. \code{getline()} for reads (tolerates blank
lines, grows with the message); \code{fwrite + '\textbackslash n' + fflush} for
writes under a per-transport mutex so a reply and a notification thread never
interleave a frame. \code{CMCP\_STDIO\_MAX\_FRAME} (default 16~MiB) caps a frame
so a peer streaming endlessly without a newline can't OOM the process --- this
protects the \emph{client} too, since it reads a forked server's stdout through
the same path.

\textbf{HTTP server} (\code{transport\_http.c}) --- a hand-rolled HTTP/1.1 server
on \code{socket()}/\code{accept()}. One \code{/mcp} endpoint: \code{POST} ->
request/response, \code{GET} with SSE -> server-push channel. It bridges the
synchronous \code{read->dispatch->write} server loop onto HTTP with a single-slot
mailbox (there is only ever one logical session per transport). This is the one
place the transport layer peeks the JSON-RPC body --- to tell a notification (202
Accepted, no body) from a request (200 + body), since HTTP forces that decision.

\textbf{HTTP client} (\code{transport\_http\_client.c}) --- the mirror, built on
\code{libcurl}. Each \code{write\_fn} is a one-shot \code{POST} with its own easy
handle (so concurrent host threads don't serialise), and a separate SSE thread
long-polls \code{GET /mcp} for server-pushed frames, pushing both onto one read
queue that \code{client.c} demultiplexes. The \code{wake\_fn} slot exists because
this transport's \code{read\_fn} parks on a condvar, not a syscall (\S10).

The HTTP parser itself lives in \code{http\_parser.c} / \code{cmcp\_http\_parser.h},
extracted as its own translation unit so a libFuzzer harness can drive it without
sockets.

\subsection{server.c + worker.c --- the server half (layer 5)}

\code{server.c} holds the tool / resource / prompt registries, the dispatch
table, and the run loop. Built-in methods: the handshake, \code{ping} (answered
\code{\{\}} \emph{before} the readiness gate so it works pre-handshake), all three
\code{*/list} + act methods, \code{resources/subscribe}/\code{unsubscribe}, and
\code{logging/setLevel}. Unknown methods -> \code{-32601}; operate-class methods
before \code{initialize} -> \code{-32600}.

Capabilities \code{tools}/\code{resources}/\code{prompts} are
\emph{auto-advertised} whenever $\ge 1$ of that kind is registered;
sub-capabilities (\code{subscribe}, \code{listChanged}, \code{logging}) stay
opt-in via \code{cmcp\_server\_set\_capabilities}. The minimal server is small:

\begin{ccode}
cmcp_server_t *s = cmcp_server_new("echo", "1.0");
cmcp_server_add_tool(s, "echo", "Echo text back",
                     "{\"type\":\"object\",\"properties\":"
                     "{\"text\":{\"type\":\"string\"}},"
                     "\"required\":[\"text\"]}",
                     echo_handler, NULL);
cmcp_server_run_stdio(s);   /* blocks: read -> dispatch -> write */
\end{ccode}

\code{worker.c} is the handler thread pool. The three handler-\emph{invoking}
methods (\code{tools/call}, \code{resources/read}, \code{prompts/get}) run on
\code{CMCP\_WORKERS} threads (default 4) so a slow handler can't stall the read
loop; everything else (handshake, lists, notifications) runs inline on the loop
thread, so the lifecycle FSM needs no lock. See \S10 for the cancellation and
timeout machinery.

\subsection{client.c + session.c --- the client half (layer 5)}

\code{client.c} is an async client with \emph{one reader thread} per
\code{cmcp\_client\_t}. Host threads call \code{cmcp\_client\_call\_async} (returns
an id) and later \code{cmcp\_client\_wait(id)}; \code{cmcp\_client\_request} is
just those two fused. The reader demultiplexes inbound frames by kind: responses
go to the matching pending entry; \code{notifications/progress} go to a per-call
subscription if its token matches, else to the generic handler; reverse requests
(\code{sampling}/\code{elicitation}/\code{roots}/\code{ping}) go to the registered
host handler or default-decline.

\code{session.c} aggregates $N$ already-handshaken clients under host-supplied
names so a multi-server host sees one flat surface. Tools are namespaced
\code{<server>:<tool>}; \code{resources/read} and \code{prompts/get} take an
explicit \code{(server, uri|name)} pair (URIs already contain colons, so folding
the server in would be ambiguous). List calls fan out async, fan in sequentially,
and follow each server's \code{nextCursor} pagination automatically.

% ---------------------------------------------------------------------
\section{The threading model}

This is the subtlest part of the codebase and the most worth studying.

\textbf{Server side.} The run loop owns the transport on one thread: read, parse,
dispatch. Dispatch forks two ways --- cheap/lifecycle methods run \emph{inline}
(so the FSM is lock-free), and the three handler-invoking methods go to a bounded
worker pool. The pool's submit \emph{blocks} when full, so backpressure
propagates to the loop instead of growing memory.

\textbf{Cooperative cancellation (a contract, not a kill).} There is deliberately
no \code{pthread\_cancel} --- async-killing a thread that may hold the writer
mutex or a \code{malloc} lock is unsafe in C. Instead an in-flight table maps
\code{request id -> ctx}; \code{notifications/cancelled} or a 200ms watchdog (past
\code{CMCP\_HANDLER\_TIMEOUT\_MS}, default 30s) only \emph{flag} a context. The
handler contract is therefore load-bearing: \textbf{a handler MUST poll
\code{cmcp\_handler\_cancelled()} in any loop or before any blocking step}. A
handler that blocks forever burns its worker slot permanently; \code{CMCP\_WORKERS}
such handlers deadlock the server.

\textbf{Client side.} One reader, many waiters. Each call gets a completion record
with its own mutex+condvar (localised wakeup, no thundering herd). Multiple async
calls may be in flight and \code{wait} may be called in any order.

\begin{diag}
  caller thread(s)                   |  reader thread
                                     |  for (;;) {
  call_async:                        |     transport_read()
    register pending id              |     rpc_parse()
    alloc + link completion record   |     switch (kind):
    send_message                     |       RESPONSE -> deliver by id, broadcast cv
    return id                        |       NOTIFICATION -> progress sub? else notif_fn
  wait(id):                          |       REQUEST -> ping/sampling/elicit/roots
    cond_wait until done             |                  else -32601
    move response out, free          |  }
\end{diag}

\textbf{Two shutdown subtleties worth internalising.} (1) On Linux, closing an fd
does \emph{not} interrupt a \code{read()} already blocked on it. So the stdio
reader is woken by a signal (\code{SIGUSR2}, non-restarting handler) that makes
the syscall return \code{EINTR}; a \code{pthread\_tryjoin\_np} retry loop
re-sends the signal to close a race window. (2) A transport whose \code{read\_fn}
parks on a \emph{userspace} primitive (the HTTP client's condvar) can't be woken
by a signal --- that's why the vtable has the optional \code{wake\_fn}, which
stdio leaves NULL.

\textbf{Callbacks run on the reader thread, serialised.} So a slow host callback
(sampling doing an LLM round-trip) stalls inbound frames until it returns. From
inside a callback the non-blocking calls (\code{call\_async}, \code{notify},
\code{cancel}) are safe, but the blocking pair (\code{request}, \code{wait})
self-deadlocks --- the thread that would complete them is the one running the
callback.

% ---------------------------------------------------------------------
\section{Memory ownership}

Manual \code{malloc}/\code{free} throughout; the rule is a single convention.

\begin{center}\small
\begin{tabular}{l l}
\toprule
Object & Who frees \\
\midrule
\code{cmcp\_json\_t*} passed into \code{*\_set}/\code{*\_request}/handlers & callee (library) \\
\code{cmcp\_json\_t*} returned from \code{*\_get} & nobody --- borrowed, lives with parent \\
\code{cmcp\_rpc\_message\_t} & holder; \code{cmcp\_rpc\_message\_clear} before scope exit \\
Tool-handler \code{*out\_content} & library, after emitting the response \\
Tool-handler \code{arguments} & caller (library) --- borrowed, handler must NOT free \\
Owned transport (\code{connect\_stdio}) & client, on \code{client\_free} \\
Spawned child process & client, reaped via SIGTERM + \code{waitpid} \\
Clients added to a session & session, on \code{session\_free} \\
\bottomrule
\end{tabular}
\end{center}

The convention: any function whose docs say ``consumed'' or ``takes ownership''
is a transfer point. Everything else is borrowed. Negative return values are
\code{cmcp\_err\_t} codes (\code{include/cmcp.h}); the library never \code{exit()}s
--- a handler must survive a single bad call without taking the process down.

% ---------------------------------------------------------------------
\section{Security and hardening}

cMCP's trust model is precise: the \emph{host} is trusted (cMCP runs in its
address space), but the \emph{peer} on the other end of a transport may be hostile
or buggy, and its input is untrusted. A STRIDE threat catalogue maps to concrete,
individually-tested controls --- each reachable from a snapshot-once environment
knob so an operator can tune or disable it without recompiling.

\begin{center}\small
\begin{tabular}{p{0.27\linewidth} p{0.46\linewidth} p{0.18\linewidth}}
\toprule
Threat & Control & Knob \\
\midrule
SSRF via discovery & egress guard on the \emph{resolved} peer; rejects
  metadata/link-local/RFC1918/CGNAT/ULA; DNS-rebinding-safe & \code{CMCP\_HTTP\_ALLOW\_PRIVATE} \\
DNS rebinding & \code{Origin} allow-list on the HTTP server & \code{CMCP\_HTTP\_ALLOWED\_ORIGINS} \\
Insecure deserialization & schema validation \emph{before} dispatch; parse depth + element caps & \code{CMCP\_JSON\_MAX\_DEPTH/\_ELEMENTS} \\
Memory bomb / OOM & stdio frame cap, HTTP body cap, handler \code{RLIMIT\_AS} & \code{CMCP\_STDIO\_MAX\_FRAME} \\
Slowloris / spike & per-recv idle timeout, cumulative deadline, accept-rate token bucket & \code{CMCP\_HTTP\_ACCEPT\_RATE} \\
Credential leakage & log redactor over credential-shaped keys & \code{CMCP\_LOG\_REDACT} \\
\bottomrule
\end{tabular}
\end{center}

The controls are layered front-to-back: a request crosses the \code{Origin}
allow-list and SSRF guard, then frame/body/depth/element caps, then schema
validation, then rate-limit and idle/deadline timeouts, before a handler runs
under \code{RLIMIT\_AS}; anything it logs passes the credential redactor on the
way out.

\textbf{Deliberate non-goals (scope, not omission).} cMCP links \emph{no} TLS
library --- a protocol library is the wrong layer to own cert rotation; put
nginx/caddy/Envoy in front (\code{docs/deployment-tls.md}). Application-layer auth
(OAuth 2.1, API keys, mTLS) is the terminator's job; cMCP trusts the
already-authenticated request. Out-of-process handler sandboxing is a separate
deferred tier --- today's threat model is single-author tools, not untrusted
third-party code. The network surface is \emph{defensive depth} for a deployment
that is stdio-primary and single-tenant.

% ---------------------------------------------------------------------
\section{How correctness is proved}

For a protocol library the central question is not ``does it run'' but ``does it
emit \emph{exactly} the bytes the spec requires, and interoperate with an
independent implementation.'' cMCP answers with cross-implementation checks in
both wire roles, not self-consistency alone.

\begin{itemize}
  \item \textbf{cMCP client vs.\ TS reference server} and \textbf{TS client vs.\
    cMCP server} (\code{make conformance}) --- both directions against the
    official TypeScript SDK.
  \item \textbf{Replay gate} (\code{make replay}) --- real wire transcripts,
    captured through the \code{cmcp-tee} proxy, replayed and diffed so a framing
    change can't regress silently.
  \item \textbf{Spec-drift watch} (\code{make check-spec-drift}, weekly CI) ---
    compares the pinned \code{CMCP\_PROTOCOL\_VERSION} against the upstream spec.
  \item \textbf{Sanitisers + fuzzers} --- the suite runs clean under Valgrind,
    ASan/UBSan, and ThreadSanitizer; four libFuzzer harnesses drive the JSON,
    JSON-RPC, schema, and HTTP-head parsers; a soak driver proves no drift over
    long stdio/HTTP runs.
\end{itemize}

The tests are not mocked: \code{test\_stdio\_roundtrip} actually \code{fork()}s a
child server and runs a real handshake; \code{test\_client\_server} wires real
\code{cmcp\_server\_run} instances onto pipe pairs to exercise the async client
and multi-server session end-to-end. The suites map onto the classic
Unit~$\rightarrow$~Integration~$\rightarrow$~Evaluation pyramid, with the security
axis cutting across all layers.

% =====================================================================
\part{A guided study path}

% ---------------------------------------------------------------------
\section{Reading order}

Read the code in the order data flows, bottom-up, with a header before each
implementation:

\begin{enumerate}
  \item \code{include/cmcp.h} --- version, protocol, the \code{cmcp\_err\_t}
    codes. Five minutes; it frames everything.
  \item \code{include/cmcp\_json.h} then \code{src/json.c} --- the value tree is
    the vocabulary every other layer speaks.
  \item \code{include/cmcp\_types.h} + \code{src/rpc.c} --- the three message
    shapes and the id table.
  \item \code{include/cmcp\_schema.h} + \code{src/schema.c} --- alongside
    \code{docs/schema-conformance.md} for the subset boundary.
  \item \code{include/cmcp\_transport.h} + \code{src/transport\_stdio.c} --- the
    vtable, on its simplest backend, before the HTTP backends.
  \item \code{include/cmcp\_server.h} + \code{src/server.c} (then
    \code{src/worker.c}) --- the registry, dispatch, and pool.
  \item \code{include/cmcp\_client.h} + \code{src/client.c} (then
    \code{include/cmcp\_session.h} + \code{src/session.c}) --- the reader-thread
    demux and multi-server aggregator.
  \item \code{examples/echo-server.c} + \code{examples/minimal-client.c} --- the
    smallest end-to-end pair; trace a \code{tools/call} through every layer above.
  \item Only then the HTTP backends (\code{transport\_http.c},
    \code{transport\_http\_client.c}) --- the most machinery, least protocol.
\end{enumerate}

Keep \code{docs/architecture.md} open as the prose ``why'' companion to the
headers' ``what.''

% ---------------------------------------------------------------------
\section{Hands-on exercises}

\begin{enumerate}
  \item \textbf{Build and self-test.} \code{make \&\& make test}. Then run one
    binary directly, e.g.\ \code{./tests/test\_schema}, and open its source to see
    the assertion style.
  \item \textbf{Watch a real handshake.} Use \code{cmcp-inspect} to spawn
    \code{filesystem-mcp}, run the handshake, and list its tools/resources. Read
    the \code{initialize} exchange it prints.
  \item \textbf{Capture the wire.} Put \code{cmcp-tee} between
    \code{cmcp-inspect} and a server; inspect the JSONL log to see every frame in
    both directions --- this is exactly what the replay gate diffs.
  \item \textbf{Write a tool.} Add a second tool to \code{examples/echo-server.c}
    with a non-trivial \code{inputSchema}, rebuild, and call it with deliberately
    bad arguments. Confirm the failure arrives as a tool-level
    \code{\{isError:true\}} result, \emph{not} a JSON-RPC error (\S2, \S9).
  \item \textbf{Break a control.} Set \code{CMCP\_JSON\_MAX\_DEPTH=4} and feed a
    deeply-nested object; watch the parser fail closed with \code{-32700}. Then
    explore the other knobs in the \S12 table.
  \item \textbf{Prove thread-safety.} Run \code{make test-tsan} and read
    \code{tests/test\_client\_server.c} to see how concurrent async calls are
    exercised.
\end{enumerate}

\textbf{Build this document:} \code{pdflatex docs/cmcp-study-guide.tex} (run twice
so the table of contents resolves). It needs no external figures.

% ---------------------------------------------------------------------
\section{Glossary}

\begin{description}[leftmargin=0pt,style=nextline,itemsep=2pt]
  \item[Host] The user-facing application that owns the LLM and decides which
    servers to trust.
  \item[Client / Server] The two endpoints of one MCP session; client drives,
    server offers. 1:1 per session.
  \item[Capability] A declared feature, exchanged at \code{initialize};
    unadvertised capabilities must not be used.
  \item[Primitive] A tool, resource, or prompt --- the three things a server
    exposes.
  \item[Tool / Resource / Prompt] Model-controlled function / application-pulled
    data / user-invoked template.
  \item[Notification] An id-less JSON-RPC message; no response expected.
  \item[Reverse request] A server-originated request to the host: sampling,
    roots, or elicitation.
  \item[Transport] The byte channel under JSON-RPC: stdio or Streamable HTTP.
  \item[Handler] User code registered for a tool/resource/prompt; runs on a
    worker thread under the cooperative-cancellation contract.
  \item[Dual error channel] Protocol failures use JSON-RPC \code{error};
    tool-run failures use an \code{isError:true} \emph{result}.
\end{description}

\vfill
{\color{rulegrey}\hrule height 0.6pt}
\vspace{4pt}
{\small\color{rulegrey}
cMCP --- pure C11, three static link targets, no third-party JSON / MCP SDK.
Companion docs: \code{docs/architecture.md} (design rationale),
\code{docs/cmcp-engineering-report.tex} (measured claims),
\code{docs/schema-conformance.md} (validator subset).}

\end{document}
